\textbf{Phần 1. Trắc nghiệm lựa chọn}
\setcounterpageref{ex}{0}

\begin{ex} %[2H9V46]
Đồ vật nào dưới đây có dạng khối cầu?
\choice
{Hộp sữa bột}
{Hộp phấn viết}
{Cục tẩy}
{\True Quả bóng đá}
\loigiai{
Quả bóng đá có dạng khối cầu.
\textbf{Đáp án của câu hỏi là D}
}
\end{ex}

\begin{ex} %[2H9V46]
Đồ vật nào dưới đây có dạng khối trụ?
\choice
{Quả địa cầu}
{\True Lon nước ngọt}
{Hộp bánh quy vuông}
{Cuốn sách giáo khoa}
\loigiai{
Lon nước ngọt có dạng khối trụ.
\textbf{Đáp án của câu hỏi là B}
}
\end{ex}

\begin{ex} %[2H9V46]
Khối cầu có đặc điểm nào dưới đây?
\choice
{Có hai mặt đáy phẳng}
{Có $6$ mặt vuông}
{\True Lăn được về mọi phía}
{Không lăn được}
\loigiai{
Khối cầu có bề mặt cong đều nên lăn được về mọi phía.
\textbf{Đáp án của câu hỏi là C}
}
\end{ex}

\begin{ex} %[2H9V46]
Khối trụ có bao nhiêu mặt phẳng hình tròn?
\choice
{\True $2$ mặt phẳng}
{$1$ mặt phẳng}
{$3$ mặt phẳng}
{Không có mặt nào}
\loigiai{
Khối trụ có hai mặt đáy là hai hình tròn phẳng bằng nhau.
\textbf{Đáp án của câu hỏi là A}
}
\end{ex}

\begin{ex} %[2H9V46]
Người ta có thể xếp chồng các khối nào lên nhau mà không bị lăn đổ?
\choice
{Các khối cầu xếp chồng lên nhau}
{Hai quả bóng xếp chồng}
{Khối cầu đặt trên khối trụ}
{\True Các khối trụ xếp chồng theo mặt đáy phẳng}
\loigiai{
Xếp chồng các khối trụ theo mặt đáy phẳng thì đứng vững chắc.
\textbf{Đáp án của câu hỏi là D}
}
\end{ex}

\begin{ex} %[2H9V46]
Trong các vật sau: quả bưởi, lon sữa ông thọ, hộp chè hình trụ, viên bi, chiếc cốc hình trụ. Có bao nhiêu vật có dạng khối cầu?
\choice
{$3$ vật}
{\True $2$ vật}
{$1$ vật}
{$4$ vật}
\loigiai{
Hai vật có dạng khối cầu là: quả bưởi và viên bi.
\textbf{Đáp án của câu hỏi là B}
}
\end{ex}

\begin{ex} %[2H9V46]
Cắt một khối trụ bằng một nhát cắt phẳng song song với hai đáy, phần nhận được ở mặt cắt có hình gì?
\choice
{Hình vuông}
{Hình tam giác}
{\True Hình tròn}
{Hình chữ nhật}
\loigiai{
Mặt cắt song song với đáy khối trụ là một hình tròn.
\textbf{Đáp án của câu hỏi là C}
}
\end{ex}

\begin{ex} %[2H9V46]
Viên phấn viết bảng thông thường có dạng hình gì?
\choice
{\True Khối trụ}
{Khối cầu}
{Khối lập phương}
{Khối hộp chữ nhật}
\loigiai{
Thân viên phấn có dạng khối trụ tròn dài.
\textbf{Đáp án của câu hỏi là A}
}
\end{ex}

\begin{ex} %[2H9V46]
Một mô hình lâu đài được ghép từ $4$ khối trụ, $3$ khối cầu và $5$ khối lập phương. Có tất cả bao nhiêu khối hình?
\choice
{$10$ khối}
{$11$ khối}
{$14$ khối}
{\True $12$ khối}
\loigiai{
$4 + 3 + 5 = 12$ (khối).
\textbf{Đáp án của câu hỏi là D}
}
\end{ex}

\begin{ex} %[2H9V46]
Vật nào sau đây vừa có thể trượt vừa có thể lăn?
\choice
{Khối cầu}
{\True Khối trụ}
{Khối lập phương}
{Khối hộp chữ nhật}
\loigiai{
Khối trụ đặt đứng thì trượt được theo đáy phẳng, đặt nằm thì lăn được.
\textbf{Đáp án của câu hỏi là B}
}
\end{ex}

\begin{ex} %[2H9V46]
Có $6$ khối trụ bằng nhau, mỗi khối cao $3\text{ cm}$. Xếp chồng cả $6$ khối trụ lên nhau thì được cột trụ cao bao nhiêu xăng-ti-mét?
\choice
{$15\text{ cm}$}
{$21\text{ cm}$}
{\True $18\text{ cm}$}
{$16\text{ cm}$}
\loigiai{
$3 \times 6 = 18\text{ (cm)}$.
\textbf{Đáp án của câu hỏi là C}
}
\end{ex}

\begin{ex} %[2H9V46]
Trái Đất nơi chúng ta đang sống có hình dạng gần giống với khối gì?
\choice
{\True Khối cầu}
{Khối trụ}
{Khối lập phương}
{Hình tròn phẳng}
\loigiai{
Trái Đất có dạng hình khối cầu.
\textbf{Đáp án của câu hỏi là A}
}
\end{ex}

\textbf{Phần 2. Trắc nghiệm đúng sai}
\setcounterpageref{ex}{0}

\begin{ex} %[2H9V46]
Xét tính đúng sai của các nhận định về khối trụ và khối cầu:
\choiceTF
{\True Khối cầu không có cạnh và không có đỉnh nào.}
{Khối trụ có thể lăn được về mọi phía giống như khối cầu.}
{\True Lon nước ngọt đặt đứng trên bàn không bị lăn vì có mặt đáy phẳng.}
{Khối cầu có thể xếp chồng lên nhau thành tháp cao vững chắc.}
\loigiai{
\begin{itemize}
    \item Ý a) Đúng: Đúng.
    \item Ý b) Sai: Khối trụ chỉ lăn được theo một hướng khi đặt nằm tròn.
    \item Ý c) Đúng: Đúng.
    \item Ý d) Sai: Khối cầu hình cong sẽ bị lăn đổ ngay.
\end{itemize}
\begin{itemize}
    \item a) \textbf{Đúng}.
    \item b) \textbf{Sai}.
    \item c) \textbf{Đúng}.
    \item d) \textbf{Sai}.
\end{itemize}
}
\end{ex}

\begin{ex} %[2H9V46]
Trong giờ học xếp hình khối, bạn Nam có một bộ đồ chơi gồm các khối hình:
\choiceTF
{\True Nam lấy 3 quả bóng tròn thì Nam đã chọn 3 khối cầu.}
{\True Nam dùng các ống bơ để làm cột trụ lâu đài thì đó là các khối trụ.}
{Xếp một khối cầu bên dưới rồi đặt khối trụ lên trên thì cột trụ đứng vững.}
{\True Một khối trụ có 2 mặt phẳng đối diện nhau là hình tròn bằng nhau.}
\loigiai{
\begin{itemize}
    \item Ý a) Đúng: Đúng.
    \item Ý b) Đúng: Đúng.
    \item Ý c) Sai: Khối cầu bên dưới tròn sẽ làm khối trụ bên trên đổ ngã.
    \item Ý d) Đúng: Đúng.
\end{itemize}
\begin{itemize}
    \item a) \textbf{Đúng}.
    \item b) \textbf{Đúng}.
    \item c) \textbf{Sai}.
    \item d) \textbf{Đúng}.
\end{itemize}
}
\end{ex}

\textbf{Phần 3. Trả lời ngắn}
\setcounterpageref{ex}{0}

\begin{ex} %[2H9V46]
Một hộp đồ chơi có 4 khối cầu và 6 khối trụ. Hỏi hộp đó có tất cả bao nhiêu khối hình?
\shortans{10}
\loigiai{
\textbf{Đáp số:} $10$\\
$4 + 6 = 10$ khối.
}
\end{ex}

\begin{ex} %[2H9V46]
Người ta xếp chồng 5 lon sữa hình trụ lên nhau, mỗi lon cao 4 cm. Hỏi cột lon sữa cao bao nhiêu xăng-ti-mét?
\shortans{20}
\loigiai{
\textbf{Đáp số:} $20$\\
$4 \times 5 = 20\text{ cm}$.
}
\end{ex}

\begin{ex} %[2H9V46]
Một chú hề tung hứng 3 quả bóng xanh và 5 quả bóng đỏ có dạng khối cầu. Hỏi chú hề tung hứng tất cả bao nhiêu quả bóng?
\shortans{8}
\loigiai{
\textbf{Đáp số:} $8$\\
$3 + 5 = 8$ quả bóng.
}
\end{ex}

\begin{ex} %[2H9V46]
Trong rổ có 15 quả cam có dạng khối cầu, mẹ lấy ra 7 quả để vắt nước. Trong rổ còn lại bao nhiêu quả cam?
\shortans{8}
\loigiai{
\textbf{Đáp số:} $8$\\
$15 - 7 = 8$ quả.
}
\end{ex}

\begin{ex} %[2H9V46]
Khối trụ có mấy mặt phẳng hình tròn?
\shortans{2}
\loigiai{
\textbf{Đáp số:} $2$\\
Có 2 mặt phẳng hình tròn ở hai đầu đáy.
}
\end{ex}

\begin{ex} %[2H9V46]
Một mô hình xe buýt đồ chơi gồm 6 bánh xe có dạng khối trụ dẹt. Hỏi 3 chiếc xe buýt như thế có tất cả bao nhiêu bánh xe?
\shortans{18}
\loigiai{
\textbf{Đáp số:} $18$\\
$6 \times 3 = 18$ bánh xe.
}
\end{ex}
